% Einheit 1 von 2 – Umkehrbarkeit, Bereiche und Term (bank/umkehrfunktion/e1.jsonl, zusammenbau v0.1)
%% TODO Zeitmarke Einheit 1: Marken-Zeile nicht lesbar
%% TODO Zweigzeile Teil 1: was hier gelernt wird (Satz in Schülersprache aus dem Einheitstitel) – Einheit 1
\einheitenkopf[e1]{Einheit 1 von 2 · Umkehrbarkeit, Bereiche und Term}
\zweigzeile{keine Prüfungsaufgabe}
\setcounter{aufgabe}{7}

%% TODO Titel: Ich-kann-Satz fehlt in der Bank (Platzhalter: Kettenname) – Nr. 8
%% TODO Anweisung: ein Satz über der ersten Teilaufgabe fehlt in der Bank – Nr. 8
\begin{aufgabe}{Umkehrbarkeit, Bereiche und Term}
%% TODO \rechenplatz{n}: Schrittzahl des Grundfalls fehlt in der Bank (nur bei fester Schreibform, 2.3 b)
\begin{teile}
\teil Die Abbildung zeigt den Graphen von f auf dem Intervall $[-2;\, 2]$. Ist f dort umkehrbar? \kreuz{umkehrbar} \\ \kreuz{nicht umkehrbar}

\begin{ksys}[xmin=-3,xmax=3,ymin=-4,ymax=4,ablesen] \funktion{0.5*\x^3}{f} \end{ksys}
\teil $f(x) = 2x - 1$. Gib den Term der Umkehrfunktion g an und berechne $g(7)$.
\teil $f(x) = e^x + 3$ auf ganz ℝ, g ist die Umkehrfunktion. Gib Definitions- und Wertebereich von g an.
\teil Gegeben ist $f(x) = 4e^{-\frac{1}{2}x} - 4$; g ist die Umkehrfunktion von f. Gib den Definitionsbereich von g an und weise nach, dass $g(x) = -2 \cdot \mathrm{ln}\left(\frac{1}{4}x + 1\right)$ gilt \hfill (Abitur 2026 LK).
\teil $f(x) = 3x \cdot e^{-\frac{1}{2}x}$ hat den Hochpunkt $H(2 | \tfrac{6}{e})$ und nähert sich für große x der x-Achse. h ist f, eingeschränkt auf $[2;\, \infty[$. Begründe, dass f nicht umkehrbar ist, h aber schon, und gib Definitions- und Wertebereich der Umkehrfunktion von h an \hfill (Abitur 2025 LK).
\end{teile}
\end{aufgabe}

%% TODO Titel: Ich-kann-Satz fehlt in der Bank (Platzhalter: Kettenname) – Nr. 9
%% TODO Anweisung: ein Satz über der ersten Teilaufgabe fehlt in der Bank – Nr. 9
\begin{aufgabe}{Radius der Flüssigkeitsoberfläche in Abhängigkeit von der Füllhöhe über die Umkehrung der Profilfunktion nachweisen}
\begin{teile}
\teil Der Längsschnitt eines Glases mit senkrechter Achse wird für $-2 \le x \le 2$ durch $p(x) = 0{,}75x^2$ beschrieben (Maße in cm). Weise nach, dass der Radius der Flüssigkeitsoberfläche bei der Füllhöhe h durch $r(h) = \frac{2}{3}\sqrt{3h}$ gegeben ist, und berechne ihn für $h = 3$ \hfill (Abitur 2017 LK).
\end{teile}
\end{aufgabe}

%% TODO Titel: Ich-kann-Satz fehlt in der Bank (Platzhalter: Kettenname) – Nr. 10
%% TODO Anweisung: ein Satz über der ersten Teilaufgabe fehlt in der Bank – Nr. 10
\begin{aufgabe}{Umkehrbarkeit, Bereiche und Term – Fehler finden, Begründen, Anwendung, Darstellungswechsel}
\begin{teile}
\teil $f(x) = e^x - 2$ auf ganz ℝ, g ist die Umkehrfunktion. Tom schreibt: „Der Definitionsbereich von g ist ℝ, wie bei f.“ Finde den Fehler.
\teil Begründe, warum eine Funktion mit einem Hochpunkt im Inneren ihres Definitionsbereichs nicht umkehrbar ist.
\teil Ein Taxi kostet 3,50 € Grundgebühr und 2,10 € je Kilometer. Gib die gefahrene Strecke als Funktion des Preises an und berechne, wie weit man für 24,50 € fährt.
\teil Die Wertetabelle gehört zu einer umkehrbaren Funktion f. Gib die Wertetabelle der Umkehrfunktion g an.

\wertetabelle[1,3,5,7]{x}{f(x)}{0,1,2,3}
\end{teile}
\end{aufgabe}
